\documentclass[aspectratio=169]{beamer}
\usetheme{metropolis}
\usepackage[utf8]{inputenc}
\usepackage[spanish]{babel}
\usepackage{booktabs}
\title{Reto 2 INE — Índice Territorial de Riesgo de Abstención}
\subtitle{Índice de riesgo de abstención por sección electoral}
\author{José Ramón García González \and Marcos Bautista Muñua \and Erick Michel Meza Mendoza \and Angel Ainar Lazcano Sanchez}
\institute{INE}
\date{31 de octubre de 2026}

\begin{document}
\maketitle

\begin{frame}{Resumen ejecutivo}
\begin{itemize}
\item 70,751 secciones electorales en 2024; 68,651 con tasa válida.
\item Modelo HistGradientBoosting con historia electoral 2009-2021 y contexto espacial municipal/estatal.
\item MAE validación territorial: 0.0323 $\pm$ 0.0018; R$^2$ = 0.784 $\pm$ 0.011.
\item Reducción del MAE frente a línea base 2021: $\approx$66.9\%.
\item 9,641 secciones en riesgo alto o muy alto.
\item Personas potencialmente movilizables: X1 = 477,315; X2 = 1,113,735.
\end{itemize}
\end{frame}

\begin{frame}{Preguntas de investigación}
\begin{enumerate}
\item ¿Qué zonas presentan mayor abstención en 2024 y qué información ayuda a priorizar?
\item ¿Qué variables históricas, sociodemográficas y territoriales explican la abstención?
\item ¿Cuántas personas podrían movilizarse con acciones X1 e X2 y dónde priorizar?
\end{enumerate}
Alcance: pronóstico retrospectivo \textit{nowcast} con información hasta 2021 para calibrar 2024.
\end{frame}

\begin{frame}{Metodología del modelo}
\begin{itemize}
\item 92 variables: historia electoral 50, tipo de sección 5, faltantes 8, contexto espacial 28, entidad.
\item Objetivo transformado con logit; predicción inversa a probabilidad.
\item HistGradientBoostingRegressor: max\_iter=300, learning\_rate=0.04, max\_leaf\_nodes=15, min\_samples\_leaf=40.
\item Validación GroupShuffleSplit por municipio, 80/20, 3 semillas.
\end{itemize}
\end{frame}

\begin{frame}{Validación y robustez}
\begin{columns}
\column{0.5\textwidth}
\begin{tabular}{lc}
\toprule
Validación & MAE \\
\midrule
Territorial 2024 & 0.0323 \\
Temporal 2018$\rightarrow$2021 & 0.0721 \\
Línea base 2021 & 0.0974 \\
\bottomrule
\end{tabular}
\column{0.5\textwidth}
\begin{itemize}
\item Cobertura intervalo 80\%: 80.0\%, ancho medio 10.1 pp.
\item Correlación ranking simple/ponderado: 0.985.
\item Error según historia: completa 0.0289, parcial 0.0361, sin 2021 0.0709.
\end{itemize}
\end{columns}
\end{frame}

\begin{frame}{Distribución del riesgo territorial}
\begin{center}
\begin{tabular}{lrrr}
\toprule
Nivel & Secciones & Lista nominal & Riesgo medio \\
\midrule
Bajo & 22,979 & 29,014,613 & 28.86\% \\
Medio & 36,031 & 51,816,586 & 41.92\% \\
Alto & 9,344 & 15,422,746 & 54.31\% \\
Muy alto & 297 & 487,761 & 69.17\% \\
Sin dato & 2,100 & 1,588,642 & — \\
\bottomrule
\end{tabular}
\end{center}
Entidades con mayor riesgo promedio: Baja California, Sonora, Chihuahua, Durango, BCS, Michoacán.
\end{frame}

\begin{frame}{Escenarios de intervención}
\begin{itemize}
\item X1 informativa: -3 puntos en alto/muy alto $\rightarrow$ 477,315 personas.
\item X2 intensiva: -7 puntos en alto/muy alto $\rightarrow$ 1,113,735 personas.
\item Abstención promedio: sin intervención 39.35\%, X1 38.93\%, X2 38.37\%.
\end{itemize}
Los escenarios son análisis de sensibilidad, no efectos causales.
\end{frame}

\begin{frame}{Recomendaciones preventivas}
\begin{enumerate}
\item Priorizar 9,641 secciones alto/muy alto por lista nominal y cobertura.
\item X1 donde la barrera sea informativa; X2 para riesgo muy alto con barreras de acceso.
\item Revisar Sinaloa, Chihuahua, Coahuila por secciones sin dato $>$5\% de lista nominal.
\item Medir antes/después con grupo de comparación.
\item Actualizar modelo con 2024 antes de proyección 2027.
\end{enumerate}
\end{frame}

\begin{frame}{Conclusiones}
\begin{itemize}
\item Modelo territorial robusto con contexto espacial; supera línea base en 66.9\%.
\item Validación temporal confirma generalización moderada a otro año.
\item Índice permite priorizar recursos con intervalos de incertidumbre explícitos.
\item Limitaciones: no causal, territorial, ahora 2024, secciones sin dato.
\end{itemize}
\end{frame}

\end{document}
